% Source: https://artofproblemsolving.com/community/c3h1702484_you_know_this
% Stable identity: aops-topic-1702484-q01
Any two distinct points determine exactly one line, and there are $\dbinom32=3$ pairs of points: $\{A,B\}$, $\{B,C\}$, $\{A,C\}$.
\begin{itemize}
\item If $A,B,C$ are \emph{not} collinear, the three lines $AB$, $BC$, $CA$ are different (if two of them coincided, that line would contain all three points). So there are $\boxed{3}$ lines, and no line passes through all three points.
\item If $A,B,C$ are collinear, all three pairs determine the same line, so there is just $\boxed{1}$ line.
\end{itemize}

\begin{asy}
usepackage("amsmath");
size(9cm);
defaultpen(fontsize(10pt));
pair A=(0,0),B=(3,0.2),C=(1.2,2);
draw(A+1.25*(A-B)--B+0.25*(B-A),linewidth(0.7));
draw(A+0.35*(A-C)--C+0.3*(C-A),linewidth(0.7));
draw(B+0.3*(B-C)--C+0.35*(C-B),linewidth(0.7));
dot(A);dot(B);dot(C);
label("$A$",A,SE);label("$B$",B,dir(70));label("$C$",C,1.5*dir(195));
label("not collinear: $3$ lines",(1.3,-1.1));
pair P=(5.5,0),Q=(6.7,0.8),R=(8.2,1.8);
draw(P+0.4*(P-R)--R+0.2*(R-P),linewidth(0.7));
dot(P);dot(Q);dot(R);
label("$A$",P,SE);label("$B$",Q,SE);label("$C$",R,SE);
label("collinear: $1$ line",(7,-1.1));
label("\mbox{}",(min(currentpicture,user=true)+max(currentpicture,user=true))/2,black); // keeps the text after the figure black when LaTeX typesets the labels inline
\end{asy}

So the book's answer $3$ assumes the points are not collinear: it counts the lines through two of the points, not lines through all three.
